\section{Discussion and limitations}
\label{sec:discussion}

\paragraph{What the framework changes.}
\method makes the transition mechanism itself an executable molecular object.
The important consequence is not validity in isolation: it is that the same
canonical successor kernel can be evaluated, constrained, interrupted,
branched, and retargeted at every non-null committed state. Birth/death and
cycle operations determine which structural changes are reachable; the learned
law determines which are plausible; the remaining-budget controller determines
which plausible continuations serve the current goal.

\paragraph{What must be learned.}
Legal support is not enough. A finite-particle controller cannot exploit a
transition with negligible base mass, and a valid rewrite system with a nearly
uniform learned law would add little chemistry. This is why successor
likelihood, analogue recovery, path overhead, operator-resolved safeguards, and
learned-versus-uniform transport precede the optimization results. The final
evidence-backed interpretation is
\resultpending{DISCUSSION_LEARNED_PROCESS}.

\paragraph{Primitive support and accelerators.}
The current basis favors small operations with direct semantics. Primitive
cycle close/open defines topology support, while ring-system restatement is a
coordinated electronic accelerator retained by a measured support audit.
Multi-attachment insertion, fragment replacement, and whole-ring macros remain
reasonable future accelerators if a frozen path-cost study demonstrates need.
Adding them changes the event timescale and edit-budget semantics, so each must
be costed and ablated rather than presented as free capacity.

\paragraph{Control is modular, not magical.}
Exact finite-horizon control is a property of an exact base kernel and exact
backward values. At molecular scale, the learned value is approximate; improved
hypervolume does not prove exact conditioning. Likewise, hard path constraints
can make an objective unreachable at a given budget. Reporting zero-value
states and lost support is part of the result, not an implementation failure to
be smoothed away.

\paragraph{Chemical scope.}
The current state is a 2D molecular graph with at most 40 active atoms.
Stereochemistry, isotope effects, radicals, explicit protonation changes, and
general multi-attachment insertion are not modeled. Valence validity and
sanitization do not imply synthesizability, conformational plausibility,
stability, or favorable pharmacology. Independent synthesis, 3D, and biological
evaluators are therefore downstream filters or objectives, not consequences of
the rewrite closure theorem.

\paragraph{Experimental scope.}
The carbon-only exact slice is designed for exhaustive algebra, not chemical
representativeness. Editing and de novo generation use separate data mixtures,
hazard contracts, checkpoints, and evaluations. External optimizers may have
different support and inductive bias; matched oracle accounting does not make
their internal processes equivalent. The final failure and sensitivity summary
is \resultpending{DISCUSSION_FAILURE_MODES}.

\section{Conclusion}
\label{sec:conclusion}

Rewrite Generator Matching treats state-dependent executable programs as the
events of a learned stochastic process. COMPOSE realizes this construction on
trans-dimensional molecular graphs and pushes equivalent event marks to a
canonical molecular-successor kernel. That kernel is the common object for
fixed-budget editing, bounded exact control, approximate dynamic steering,
Pareto branching, and pathwise constraints. The empirical conclusion of the
frozen study is \resultpending{FINAL_EMPIRICAL_CONCLUSION}. More broadly, the
framework suggests that when a domain already provides executable,
constraint-preserving rewrites, those rewrites can serve as the native event
language of generation rather than as a post-hoc repair mechanism.
